\chapter{Constant mean curvature and constrained area}\label{sec:cmc-variation}
Finite-perimeter sets retain the enclosed region as well as its boundary. This is essential for the area-minus-volume functional, and connects the BV foundations with the CMC min-max construction of Zhou--Zhu \cite{ZZCMC}.

\begin{definition}[The area-minus-volume functional]\label{def:cmc-functional}
For a finite-perimeter set $\Omega$ in a smooth closed $(M^{n+1},g)$ and a prescribed number $c$, define
\[
 \mathcal A^c(\Omega)=P(\Omega;M)-c\operatorname{Vol}(\Omega).
\]
On a smooth boundary $\Sigma=\partial\Omega$ we use the outward unit normal and $H=\operatorname{div}_\Sigma\nu$. A CMC hypersurface has constant scalar $H$. Prescribing $c$ and finding a critical point of $\mathcal A^c$ is a different problem from minimizing perimeter at a fixed volume. See \cite{ZZCMC}, Introduction and Section 1.
\end{definition}

\begin{proposition}[The smooth CMC Euler-Lagrange equation]\label{prop:cmc-first-variation}
For a smooth compact region $\Omega\subset M$ with smooth boundary, a variation with boundary velocity $f\nu$ satisfies
\[
 \delta\mathcal A^c(f\nu)=\int_{\partial\Omega}(H-c)f\,dA.
\]
It is critical under all smooth variations if and only if $H=c$ everywhere.
\end{proposition}
\begin{proof}
Proposition \ref{prop:smooth-first-variation} gives the area derivative $\int Hf$. The divergence theorem applied to an extension $X$ of $f\nu$ gives the volume derivative
$\int_\Omega\operatorname{div}_M X=\int_{\partial\Omega}f$.
Subtract $c$ times this derivative. Every smooth $f$ extends to a variation, so vanishing for all $f$ is equivalent to $H-c=0$.
\end{proof}

\begin{definition}[Almost embedded hypersurfaces]\label{def:almost-embedded}
An almost embedded hypersurface is a smooth immersion whose local sheets are embedded and, when they meet, lie on one side of each other: tangential touching is allowed but transverse crossing is excluded. This is the notion used in \cite{ZZCMC}, following Theorem 0.1. It is weaker than embeddedness. The multiplicity of a parametrized sheet and the density at a touching point should not be confused.
\end{definition}

\begin{theorem}[Prescribed positive constant mean curvature]\label{thm:zhou-zhu-cmc}
Let $M^{n+1}$ be a smooth closed connected Riemannian manifold with $2\leq n\leq6$. For each $c>0$ there exists a nonempty smooth closed almost embedded hypersurface with constant mean curvature $c$, in the boundary orientation of the Zhou--Zhu construction. The min-max hypersurface has sheet multiplicity one. Source: \cite{ZZCMC}, Theorem 0.1 and Remark 0.2, in the supplied arXiv v2 of August 11, 2017. Proof external. The conclusion permits touching; it does not assert global embeddedness or connectedness.
\end{theorem}

\chapter{Constant geodesic curvature networks}\label{sec:curvature-networks}
The one-dimensional problem has a separate regularity theory. In particular, a hypersurface theorem in ambient dimensions three through seven cannot simply be applied with $n=1$.

\begin{theorem}[Curvature networks on a closed surface]\label{thm:zhou-zhu-networks}
Let $(S,g)$ be a smooth closed connected Riemannian surface and let $c>0$. There is a nontrivial integral one-varifold induced with multiplicity one by the boundary of a finite-perimeter set, and a finite set $F\subset S$, such that away from $F$ its boundary is an almost embedded curve of constant geodesic curvature $c$. At each point of $F$ the density is an integer and every tangent cone is a stationary geodesic network in the tangent plane, smooth away from the origin. Source: \cite{ZZNet}, Theorem 1.1. Its proof is external. This result does not assert that the entire output is a smooth embedded closed curve.
\end{theorem}
